\section{综合结果分析与业务建议}

本章的目的在于把分类、回归和聚类三类任务放在同一分析框架中解释，并给出谨慎的行为管理建议。所有建议都基于模型结果和描述性统计，不使用因果化表述，也不将模型输出用于医学诊断。

\subsection{三类任务的共同信息}

分类任务显示，在排除结果变量泄漏后，设备使用时长、睡眠时长和相关工程特征仍能为高风险筛查提供一定信息。最终 Gradient Boosting 在 threshold=0.14 下取得 Recall=0.6420 和 F1=0.5355，说明模型可以识别一部分高风险样本，但仍存在 63 个漏判和 133 个误报。

回归任务显示，不同目标变量的可预测性差异很大。\texttt{digital\_dependence\_score} 的最佳 $R^2=0.9839$，说明该目标在当前合成数据中与设备使用、通知和解锁变量高度一致；而 \texttt{productivity\_score} 的最佳 $R^2=-0.0041$，说明生产力得分不能被当前特征稳定预测。这个对比提醒我们，不能因为某个目标预测效果好，就推断所有数字生活方式结果都能被有效预测。

聚类任务给出了三类探索性画像：高社交媒体使用型、高设备依赖型和低负荷均衡型。但 KMeans $k=3$ 的 Silhouette=0.1860，聚类结构偏弱，因此画像更适合用于描述和讨论，而不是用于严格分群或自动化决策。

\subsection{数字生活方式管理建议}

从描述性分析和模型结果看，设备使用时长、通知暴露、解锁频率、社交媒体时长、睡眠时长和睡眠质量是报告中反复出现的变量。基于这些发现，可以提出以下温和建议：

\begin{itemize}
  \item 关注设备使用时长。设备使用时长在分类和数字依赖回归中都具有较高解释价值，因此用户可以通过记录日均使用时间来了解自身数字负荷。
  \item 管理通知干扰。通知数量与数字依赖得分预测关系较强，减少不必要通知可能有助于降低注意力被打断的频率，但本文不声称通知减少必然导致健康改善。
  \item 合理安排社交媒体使用。聚类中的高社交媒体使用型具有明显的行为特征，说明社交媒体时长适合作为用户画像维度之一。
  \item 保证睡眠和活动。EDA 显示高风险组与非高风险组在睡眠时长和睡眠质量上存在描述性差异，因此睡眠和活动变量应在数字生活方式管理中被共同考虑。
  \item 对筛查结果进行人工复核。分类模型存在误报和漏判，模型输出只能作为初筛参考，不能替代人工判断、专业评估或个体实际情况。
\end{itemize}

\subsection{方法局限}

本文主要存在三方面限制。第一，数据为合成数据，变量关系可能反映生成机制而非真实世界机制。第二，本文使用传统机器学习模型和表格特征，未引入更复杂的时间序列或个体追踪信息，因此无法刻画长期变化。第三，聚类轮廓系数较低，用户画像解释需要保持谨慎。后续若有真实纵向数据，可以进一步检验模型在真实场景中的泛化能力。
